% Source: https://artofproblemsolving.com/community/c3h2171664_the_lowest_value
% Stable identity: aops-topic-2171664-q01
The equation gives $y=\dfrac{x(x-4)}{8}$, which must be a positive integer, so $x>4$ and $8\mid x(x-4)$.
\begin{itemize}
\item If $x$ is odd, $x(x-4)$ is odd: impossible.
\item If $x=2k$, then $x(x-4)=4k(k-2)$, and $8\mid4k(k-2)$ means $k(k-2)$ is even, so $k$ is even and $4\mid x$.
\end{itemize}
So $x$ is a multiple of $4$ but not of $8$: $x\in\{12,20,28,\ldots\}$ (and $x>4$). For $x=12$ we get $y=\dfrac{12\cdot8}{8}=12=x$, which is not allowed because $x$ and $y$ must be different. For $x=20$ we get $y=\dfrac{20\cdot16}{8}=40\ne20$. So the smallest $x$ is $20$, and the smallest value of $\sqrt x$ is $\sqrt{20}=\boxed{2\sqrt5}$.
